% =============================================================
%  综合实践：两次必须动手的活动
% =============================================================
\Chapter{隐藏关\quad 画地图，花真钱}
\noindent{\small\kaishu\color{tipc}（课本两次综合实践：画校园路线图 P53--P55；欢乐购物 P94--P97）}

\begin{blunder}
带小豆角去商场，找厕所绕了三圈。最后是小豆角指着一块指示牌说：
“爸爸，刚才爸爸和小豆角是从那边过来的。”
——一个小孩的方向感比爸爸好，这事豆角爸爸认了。
\end{blunder}

\begin{gamble}{课前协商：藏宝图}
\begin{itemize}
  \item 小豆角画一张家里的“藏宝图”，标上东南西北，把一样他自己挑的东西藏起来。
  \item \textbf{豆角爸爸只能看小豆角的图去找}，不许开口问。
  \item 找到了归谁、找不到归谁，写在下面；但要小豆角说清“图上哪儿画错了”。
\end{itemize}
\negolines
\end{gamble}

\begin{bonus}
\begin{itemize}
  \item \textbf{跳蚤市场}：把家里闲置的小玩具、旧绘本摆摊卖掉，
        \textbf{赚的钱全归小豆角}。定价、算钱、找零，全小豆角自己来。
  \item \textbf{给爸爸画一张“到暖暖家怎么走”}：不许写地址，只许画图加方向。
        爸爸真按图走一次，能不能找到门？
  \item \textbf{超市寻宝}：拿手机指南针，只给三条线索（“往北走三个货架”“在收银台西边”），
        让小豆角自己找一件指定的东西。
\end{itemize}
\end{bonus}

\begin{unitinfo}
课本里安排了两次综合实践，一共占 5 页。这不是装饰——
\textbf{期末卷里的“方向”和“人民币”，几乎全从这两课出}。\\
\textbf{活动一}：画一张能用的路线图（学完第四单元后做，40 分钟）\\
\textbf{活动二}：真刀真枪买东西（学完第八单元后做，60 分钟）
\end{unitinfo}

\subsection{活动一：认方向，然后画一张真能用的图}

\begin{lesson}{1}{40}
\textbf{今天要拿下的}：认得出东南西北；画的图别人拿着能找得到地方。

\begin{materials}
手机（开指南针）、A3 纸或两张 A4 拼起来、铅笔、直尺。
\end{materials}

\begin{enumerate}
  \item \textbf{先在屋里认四个方向}：
        \begin{itemize}
          \item 打开手机指南针，找到\textbf{北}。让小豆角站在客厅中间，面向北。
          \item 教一句口诀：\textbf{面朝北，背后是南，左手是西，右手是东}。
          \item 让小豆角原地转圈，每转 90° 说一个方向。
          \item \textbf{考小豆角四道}：“小豆角面向南，左手是哪边？”（东）
        \end{itemize}
  \item \textbf{到小区里实地走 15 分钟}：从家门口出发，往北走 30 米看看有什么，
        再往东走看看有什么。\textbf{边走边在纸上记}：
        “北：小卖部；东：篮球场；南：大门；西：花坛”。
        \textbf{必须真走，不能坐桌边编。}
  \item \textbf{画图（重点）}：
        \begin{enumerate}[label=第\arabic*步：]
          \item 在纸的\textbf{右上角}画一个“北”箭头（\tikz[baseline=-0.6ex]{\draw[tipc,thick,->](0,0)--(0,0.5);} 朝上）。
                \textbf{这一步不画，整张图作废，考试也要扣分。}
          \item 定一个中心点（家／大门），画它四周的建筑。
          \item 用箭头把“从门口到目标地点”的路线画出来。
          \item 目标地点打一个 \textbf{▲}，写上名字。
        \end{enumerate}
        两条规则：\textbf{上北下南、左西右东}。
  \item \textbf{验收}：让家里另一个人只看这张图去找那个地方。
        \textbf{找到了，图才算合格}；找不到，问小豆角“小豆角哪儿没看懂”，回去改。
  \item \textbf{补中间方向}（课本 P55 的问题）：“种植园在西和南之间，画在哪？”
        （\textbf{西南角}）让小豆角把西南、东南、东北、西北四个角都标出来。
\end{enumerate}

\begin{saybox}
“小豆角面向北，左手是哪边？”\\
“为什么要把‘北’画在上面？”\\
“其他人拿着这张图能找到地方吗？小豆角哪儿没看懂？”
\end{saybox}

\begin{passbox}
\begin{itemize}
  \item 能说出“面朝北，左手西，右手东”，并推得出别的朝向。
  \item 画出的图有方向标、有建筑、有箭头，别人照着能找到。
\end{itemize}
\end{passbox}
\end{lesson}

\subsection{活动二：真刀真枪买东西}

\begin{lesson}{2}{60}
\textbf{今天要拿下的}：认得钱、会换钱、会算找零，并且在真买卖里把乘除加减用一遍。

\begin{materials}
附录 C 的模拟人民币（卡纸印好剪开）、家里小物件 8--10 件
（玩具、本子、橡皮、零食）、价格标签纸、笔、两张记录表。
\end{materials}

\begin{enumerate}
  \item \textbf{认钱（10 分钟）}：
        \begin{itemize}
          \item 把剪好的钱摊开，让小豆角说出每一张的面值和单位。
          \item \textbf{记住换算}：\ $1\ \text{元}=10\ \text{角}$，\ $1\ \text{角}=10\ \text{分}$。
          \item \textbf{换钱五道}（口答）：1 元换几张 5 角？10 元换几张 1 元？
                5 元换几张 5 角？2 元 5 角是多少角？38 角是几元几角？
                答案：2、10、10、25、3 元 8 角。
        \end{itemize}
  \item \textbf{给东西定价（10 分钟）}：让小豆角自己给每件东西定价，
        \textbf{要求是整元或“X 元 5 角”}。定价时想三件事：新旧、原来多少钱、
        小朋友买得起吗。
        \begin{center}
        \small
        \begin{tabular}{@{}p{4cm}p{3cm}p{3cm}p{3cm}@{}}
        \toprule
        \textbf{豆角爸爸卖的} & \textbf{原价} & \textbf{豆角爸爸定的价} & \textbf{卖掉了吗} \\
        \midrule
        \blank[3.2cm] & \blank[2.2cm] & \blank[2.2cm] & \blank[2.2cm] \\
        \blank[3.2cm] & \blank[2.2cm] & \blank[2.2cm] & \blank[2.2cm] \\
        \blank[3.2cm] & \blank[2.2cm] & \blank[2.2cm] & \blank[2.2cm] \\
        \blank[3.2cm] & \blank[2.2cm] & \blank[2.2cm] & \blank[2.2cm] \\
        \bottomrule
        \end{tabular}
        \end{center}
        \textbf{让小豆角自己体会}：“定价太高没人买，太低小豆角会亏。”
  \item \textbf{准备零钱（5 分钟）}：让小豆角自己决定每种面值准备几张，
        并算出总额。顺便说出课本那句话：“\textbf{要多准备小面值的，方便找钱}。”
  \item \textbf{正式买卖（20 分钟）——必须真给钱、真找钱}：
        \begin{itemize}
          \item 第一轮：\textbf{小豆角当卖家}，豆角爸爸当买家。豆角爸爸给大票
                （买 7 块的给 10 块），让小豆角算找零。\letgo{豆角爸爸偶尔会把找零算错}，
                看小豆角能不能发现。
          \item 第二轮：\textbf{小豆角当买家}，给小豆角 20 元预算，买三样东西，
                自己算一共多少钱、够不够、该找多少。
        \end{itemize}
        \begin{center}
        \small
        \begin{tabular}{@{}p{2.8cm}p{1.8cm}p{1.8cm}p{1.8cm}p{1.8cm}@{}}
        \toprule
        \textbf{东西} & \textbf{定价} & \textbf{付了} & \textbf{找零} & \textbf{对不对} \\
        \midrule
        \blank[2.4cm] & \blank[1.2cm] & \blank[1.2cm] & \blank[1.2cm] & \blank[1.2cm] \\
        \blank[2.4cm] & \blank[1.2cm] & \blank[1.2cm] & \blank[1.2cm] & \blank[1.2cm] \\
        \blank[2.4cm] & \blank[1.2cm] & \blank[1.2cm] & \blank[1.2cm] & \blank[1.2cm] \\
        \bottomrule
        \end{tabular}
        \end{center}
  \item \textbf{算总账（10 分钟）}：卖出所得加起来，看赚了多少；买花掉的加起来，
        看剩下多少。\textbf{顺手把两位数加减复习一遍。}
  \item \textbf{真去一趟超市（这周内）}：给小豆角 20 元，让小豆角自己买一次东西，
        \textbf{自己算找零}。\textbf{这是本单元最值钱的一次作业}，比十张卷子有用。
  \item \textbf{加点料}（课本 P97）：人民币到现在发行了五套，
        第五套纸币背面分别是西湖三潭印月、泰山、长江三峡、桂林山水、
        布达拉宫、人民大会堂。\textbf{真拿几张出来看看背面。}
\end{enumerate}

\begin{saybox}
“1 元等于几角？那 2 元 5 角呢？”\\
“豆角爸爸给小豆角 10 元，东西 7 元，小豆角该找多少？”\\
“找的零钱对不对？用加法验一下。”（找零 $+$ 货款 $=$ 付出的钱）\\
“定价 60 元，会有人买吗？”
\end{saybox}

\begin{passbox}
\begin{itemize}
  \item 五道换钱全对（包括“38 角 = 3 元 8 角”）。
  \item 能独立完成一次买卖：算总额、算找零、验算对。
  \item 能用 20 元买三样东西并算清剩多少。
\end{itemize}
\end{passbox}
\end{lesson}

\begin{tipbox}
\textbf{两个提醒}：\\
① \textbf{一定要用真钱或仿真币，别在纸上“假装”}。二年级的钱感是摸出来的。\\
② 找零算错时\textbf{别直接说答案}，说“重新数一遍钱”。
把钞票一张张摊开数，比讲道理快十倍。
\end{tipbox}
